% ── sections/02_methods.tex ───────────────────────────────────────────────────
% Concise Methods for NeurIPS two-column format

\section{Methods}

\subsection{Datasets and preprocessing}

We use six PPI networks from the EMOGI benchmark\cite{emogi2021}: CPDB (13,627 nodes),
IRefIndex (17,013), IRefIndex-2015 (12,129), Multinet (14,398), PCNet (19,781), and
STRINGdb (13,179). Node features are 64-dimensional multi-omics vectors: mutation
frequency (MF), DNA methylation (METH), gene expression (GE), and copy-number alteration
(CNA), each across 16 TCGA cancer types. Features are reordered to a common
omics-by-cancer-type schema across networks. No additional feature scaling or variance
filtering was applied in the experiments reported here. The last network in the dataset
list supplies the held-out test labels; test genes are removed from the pooled training
set, and 10\% of the remaining eligible genes are assigned to validation.

\subsection{EMGNN architecture}

The two-stage EMGNN architecture\cite{emgnn2023} processes each PPI network through a
shared $L$-layer GNN encoder (GCN default\cite{gcn2017}; GAT\cite{gat2018},
GIN\cite{gin2019}, GraphSAGE\cite{sage2017} also supported). After the input projection,
node embeddings are scaled by learnable network weights
$\mathbf{w}=\mathrm{softmax}(\boldsymbol{\theta})$ before the shared per-network message
passing stack. The resulting embeddings are connected to one meta-node per unique gene.
A single meta-graph convolution combines input-node and meta-node information, followed
by a linear classifier and log-softmax output (Fig.~1).

\subsection{Architectural improvements}

Residual skip connections\cite{resnet2016} follow all GNN layers after the first.
Normalisation is controlled by \texttt{norm\_type} (batch, graph\cite{graphnorm2021},
layer, or none). Label smoothing ($\varepsilon = 0.05$) and cosine
annealing\cite{cosine2017} ($\eta_0 = 5\times10^{-3}$ to $\eta_{\min} = 10^{-5}$) are
applied. Gradients are clipped to $\|\nabla\|_2 \le 1.0$. Adam\cite{kingma2014adam}
(weight decay $5\times10^{-4}$) trains for 2,000 epochs with early stopping
(patience~250). Default: hidden channels~=~64, $L$~=~3, dropout~=~0.5. Bayesian
hyperparameter optimisation uses Optuna (TPE sampler, median pruner, 50 trials).

\subsection{Interpretability}

Integrated Gradients\cite{ig2017} attributes predictions to input features using a
zero-vector baseline with 50 Riemann steps, aggregated across cancer meta-nodes by mean
absolute value. Gene Set Enrichment Analysis uses Enrichr ORA\cite{enrichr2013} (top 200
predicted genes) and pre-ranked GSEA\cite{gsea2005} against MSigDB Hallmark 2020\cite{msigdb2015}
(FDR~$<$~0.05).

\subsection{Advanced GNN ablation}

Nine techniques are modularly implemented as extensions to EMGNNImproved.
Five are evaluated in the ablation study:
(1) Heterophily-aware gating\cite{sgcd2024} fusing low-pass and high-pass filtered
signals via a learned gate;
(2) Cross-network attention scoring each gene's embedding against a learned network
embedding via scatter-softmax;
(3) DropEdge\cite{dropedge2020} ($p=0.1$) for structural regularisation;
(4) GraphMAE\cite{graphmae2022} self-supervised pretraining (50\% mask ratio, 200
pretraining epochs, scaled cosine error);
(5) Focal Loss\cite{focal2017} ($\gamma=2$, $\alpha=0.75$) for class imbalance.
Each is tested against the six-network baseline (GCN, norm=none, residual=True,
cosine LR, label smoothing 0.05) with 3 seeds (72, 1, 2); baseline robustness uses
5 seeds (adding 42, 99). These runs predate the fixed split-manifest protocol and are
therefore treated as legacy evidence in the final scoped evaluation. Fixed-split
reruns are supported by the released code but are not part of the present experiment
scope.

Four additional techniques are implemented but await external data or debugging:
(6) Pathway hypergraph encoding\cite{hgnn2019} from MSigDB GMT files via two-stage
incidence-matrix message passing;
(7) HIPGNN-inspired\cite{hipgnn2025} anomaly detection with spectral and spatial views;
(8) PINNACLE pretrained protein embeddings\cite{pinnacle2024} (128-dim);
(9) GNNExplainer\cite{gnnexplainer2019} for edge-level interpretability.

\subsection{Hardware}

NVIDIA RTX 5090 (32~GB), CUDA 12.8, PyTorch 2.7, PyG 2.7, Captum 0.8.0\cite{captum2020},
GSEApy 1.1.13. Code: \url{https://github.com/EthyleneC2H4/NUS-Capstone}.

% ── Architecture figure (full-width) ─────────────────────────────────────────
\begin{figure}[htbp]
\centering
\begin{tikzpicture}[
  >=Stealth,
  box/.style={draw, rounded corners=2pt, minimum height=0.45cm, text centered, font=\scriptsize},
  ppi/.style={box, fill=blue!8, minimum width=1.2cm},
  enc/.style={box, fill=green!8, minimum width=1.2cm},
  meta/.style={box, fill=violet!8, minimum width=4.5cm},
  cls/.style={box, fill=red!8, minimum width=3.5cm},
  arr/.style={->, thick, gray!60},
]
\node[ppi] (g1) at (0, 0) {$\mathcal{G}^{(1)}$};
\node[ppi] (g2) at (2, 0) {$\mathcal{G}^{(2)}$};
\node[font=\scriptsize] at (3.3, 0) {$\cdots$};
\node[ppi] (gk) at (4.8, 0) {$\mathcal{G}^{(K)}$};
\node[enc] (e1) at (0, -1.0) {GCN};
\node[enc] (e2) at (2, -1.0) {GCN};
\node[enc] (ek) at (4.8, -1.0) {GCN};
\draw[decorate, decoration={brace, amplitude=3pt, mirror}, thick, blue!50]
  ([yshift=-2pt]e1.south west) -- ([yshift=-2pt]ek.south east)
  node[midway, below=3pt, font=\tiny\itshape, text=blue!60] {shared, $L$ layers};
\node[font=\tiny, text=orange!70!black] (w1) at (0, -2.0)
  {$w_1\mathbf{H}^{(1)}$};
\node[font=\tiny, text=orange!70!black] (w2) at (2, -2.0)
  {$w_2\mathbf{H}^{(2)}$};
\node[font=\tiny, text=orange!70!black] (wk) at (4.8, -2.0)
  {$w_K\mathbf{H}^{(K)}$};
\node[meta] (mg) at (2.4, -2.8)
  {Meta-Graph GCN};
\node[cls] (mlp) at (2.4, -3.6)
  {MLP $\to$ $P(\mathrm{cancer}\mid g_i)$};
\draw[arr] (g1) -- (e1);
\draw[arr] (g2) -- (e2);
\draw[arr] (gk) -- (ek);
\draw[arr] (e1) -- (w1);
\draw[arr] (e2) -- (w2);
\draw[arr] (ek) -- (wk);
\draw[arr] (0, -2.1) -- ([xshift=-1.2cm]mg.north);
\draw[arr] (2, -2.1) -- (mg.north);
\draw[arr] (4.8, -2.1) -- ([xshift=1.2cm]mg.north);
\draw[arr] (mg) -- (mlp);
\node[font=\tiny\bfseries, text=gray!50, rotate=90, anchor=south] at (-0.8, -1.0) {Stage 1};
\node[font=\tiny\bfseries, text=gray!50, rotate=90, anchor=south] at (-0.8, -3.2) {Stage 2};
\end{tikzpicture}
\caption{\textbf{EMGNNImproved architecture.} Stage 1: $K$ PPI networks processed
by a shared GCN encoder; embeddings scaled by learnable weights $w_k$. Stage 2:
weighted embeddings aggregated into a meta-graph, processed by a second GCN, and
classified.}
\label{fig:architecture}
\end{figure}
